\clearpage
\section{연습문제}
\label{sec:construction-exercises}

문제는 A, B, C의 세 단계로 나뉜다. 별표 문제는 뒤의 두 절에 상세 풀이가 있다.

\subsection*{A. 계산과 기본 판정}

\begin{exercise}[*]
$0,1$에서 출발하여 다음 길이가 작도가능함을 체의 사칙연산과 제곱근 연산으로 설명하라.
\[
  \frac37,\qquad \sqrt3,\qquad
  \frac{1+\sqrt5}{2},\qquad
  \sqrt{2+\sqrt3}.
\]
\end{exercise}

\begin{exercise}
중심이 $a\in F\subseteq\C$이고 반지름이 $|u-v|$인 원과, $b,c\in F$를 지나는 직선의 교점 매개변수 $t$가 이차방정식을 만족함을 직접 전개하라. $F$는 복소켤레에 닫혀 있다고 가정한다.
\end{exercise}

\begin{exercise}[*]
$\sqrt[3]2$가 자와 컴퍼스로 작도가능하지 않음을 증명하라. 같은 논증이 $\sqrt[3]{16}$에는 왜 그대로 적용되지 않는지도 설명하라.
\end{exercise}

\begin{exercise}
$\sqrt[12]2$와 $\sqrt[16]2$의 작도가능성을 각각 판정하라. 필요한 최소다항식의 기약성도 확인하라.
\end{exercise}

\begin{exercise}[*]
삼중각 공식을 이용하여 $y=2\cos20^\circ$의 최소다항식을 구하고, $60^\circ$가 자와 컴퍼스로 삼등분되지 않음을 증명하라.
\end{exercise}

\begin{exercise}
$90^\circ$는 자와 컴퍼스로 삼등분할 수 있다. 이 사실이 문제 \thechapter.5의 결론과 모순되지 않는 이유를 설명하라.
\end{exercise}

\begin{exercise}[*]
$\zeta=\zeta_5$와 $x=\zeta+\zeta^{-1}$에 대해
\[
  x^2+x-1=0
\]
을 유도하고
\[
  \cos72^\circ=\frac{\sqrt5-1}{4},
  \qquad
  \cos36^\circ=\frac{\sqrt5+1}{4}
\]
를 증명하라.
\end{exercise}

\begin{exercise}
다음 정다각형의 작도가능성을 오일러 함수로 판정하라.
\[
  n=18,20,24,30,34,45,51,85.
\]
\end{exercise}

\begin{exercise}[*]
$X^4-2$의 분해체와 갈루아군을 구하고, 모든 근이 작도가능함을 두 가지 방법으로 설명하라.
\end{exercise}

\begin{exercise}
$f=X^4-X-1$에 대해 다음을 보이라.
\begin{enumerate}[label=(\alph*)]
  \item $f$는 $\Q$ 위에서 기약이다.
  \item 삼차분해식은 $Z^3+4Z+1$이고 기약이다.
  \item $\Disc(f)=-283$이다.
  \item $\Gal(f/\Q)\cong S_4$이고 근들은 작도가능하지 않다.
\end{enumerate}
\end{exercise}

\subsection*{B. 핵심 정리의 증명}

\begin{exercise}[*]
$\mathcal C(M)$이 덧셈, 뺄셈, 곱셈과 역수에 닫혀 있음을 기하학적으로 설명하라. 곱셈과 역수에서는 닮은삼각형을 사용하라.
\end{exercise}

\begin{exercise}
두 원의 방정식을 빼서 공통현의 직선방정식을 얻고, 두 원의 교점 계산이 직선--원 교점 계산으로 환원됨을 증명하라.
\end{exercise}

\begin{exercise}[*]
작도가능성의 기본정리에서
\[
  z\in\mathcal C(M)
  \quad\Longleftrightarrow\quad
  z\text{가 유한 이차확장 탑 안에 놓임}
\]
을 증명하라. 양방향에서 필요한 기하학적·대수적 사실을 분명히 적어라.
\end{exercise}

\begin{exercise}
유한 이차확장 탑의 정상폐포 차수가 $2$의 거듭제곱임을 탑의 길이에 대한 귀납법으로 증명하라.
\end{exercise}

\begin{exercise}[*]
유한 $2$-군 $G$가
\[
  G=G_0\trianglerighteq G_1\trianglerighteq\cdots
  \trianglerighteq G_r=1,
  \qquad [G_i:G_{i+1}]=2
\]
인 정규열을 가짐을 증명하라. 이를 이용해 유한 갈루아 $2$-확장이 이차확장 탑으로 분해됨을 보여라.
\end{exercise}

\begin{exercise}
``대수적 수의 차수가 $2$의 거듭제곱이면 작도가능하다''라는 명제가 거짓임을 $X^4-X-1$을 사용하여 설명하라. 올바른 충분조건을 적어라.
\end{exercise}

\begin{exercise}[*]
$f\in\Q[X]$가 분리인 기약 사차다항식이라 하자. 그 근이 작도가능할 필요충분조건이
\[
  \Gal(f/\Q)\in\{V_4,C_4,D_4\}
\]
임을 증명하라.
\end{exercise}

\begin{exercise}
작도가능한 대수적 수 $\alpha$의 모든 $\Q$-켤레가 작도가능함을 증명하라. 이 명제가 단순히 $[\Q(\alpha):\Q]$만으로는 보이지 않는 이유를 설명하라.
\end{exercise}

\begin{exercise}[*]
정$n$각형이 작도가능할 필요충분조건이
\[
  \varphi(n)=2^r
\]
임을 원분체 $\Q(\zeta_n)$와 갈루아 대응을 사용하여 증명하라.
\end{exercise}

\begin{exercise}
가우스--방첼 정리를 증명하라. 즉 $\varphi(n)$이 $2$의 거듭제곱일 필요충분조건이
\[
  n=2^m p_1\cdots p_r
\]
이고 $p_i$들이 서로 다른 페르마 소수인 것임을 보여라.
\end{exercise}

\subsection*{C. 구조와 종합}

\begin{exercise}[*]
표준 작도가능체 $\mathcal C_0$가 $\Q$의 무한 갈루아 확장이고, 유한 갈루아 $2$-확장들의 합집합임을 증명하라.
\end{exercise}

\begin{exercise}
$E/\Q$와 $F/\Q$가 유한 갈루아 $2$-확장이면 합성체 $EF/\Q$와 교집합 $E\cap F/\Q$도 갈루아 $2$-확장임을 증명하라.
\end{exercise}

\begin{exercise}[*]
정$m$각형과 정$n$각형이 모두 작도가능하다고 하자. 정$\gcd(m,n)$각형과 정$\operatorname{lcm}(m,n)$각형도 작도가능함을 두 방법으로 증명하라.
\begin{enumerate}[label=(\alph*)]
  \item 가우스--방첼 소인수분해를 사용한다.
  \item 원분체와 합성체를 사용한다.
\end{enumerate}
\end{exercise}

\begin{exercise}
$u=2\cos(2\pi/7)$가
\[
  u^3+u^2-2u-1=0
\]
을 만족함을 단위근 관계에서 유도하라. 이 다항식이 기약임을 보이고 정칠각형이 작도 불가능함을 결론내려라.
\end{exercise}

\begin{exercise}[*]
정구각형의 작도 불가능성과 $60^\circ$의 삼등분 불가능성이 서로 동치에 가까운 이유를 설명하라. $2\cos(2\pi/9)$의 최소다항식도 구하라.
\end{exercise}

\begin{exercise}
정삼각형과 정오각형의 작도를 결합하여 정십오각형을 작도할 수 있음을 보이라. 중심각의 정수 선형결합을 구체적으로 찾아라.
\end{exercise}

\begin{exercise}[*]
$f\in\Q[X]$가 분리인 기약 사차다항식이고 삼차분해식이 기약이며 $\Disc(f)$가 $\Q$의 제곱이라고 하자. $\Gal(f/\Q)\cong A_4$임을 보이고, $f$의 근이 작도가능하지 않음을 증명하라.
\end{exercise}

\begin{exercise}
대수적 수 $\alpha$에 대해 다음이 동치임을 증명하라.
\begin{enumerate}[label=(\roman*)]
  \item $\alpha$가 작도가능하다.
  \item $\alpha$의 최소다항식의 분해체 차수가 $2$의 거듭제곱이다.
  \item $\alpha$의 모든 켤레를 포함하는 어떤 이차확장 탑이 존재한다.
\end{enumerate}
\end{exercise}

\begin{exercise}[*]
다음 수 또는 점의 작도가능성을 판정하고 이유를 적어라.
\begin{enumerate}[label=(\alph*)]
  \item $\sqrt2+\sqrt3$,
  \item $\sqrt[3]2$,
  \item $2\cos(2\pi/17)$,
  \item $X^4-X-1$의 한 근,
  \item $\sqrt\pi$.
\end{enumerate}
\end{exercise}

\begin{exercise}
$E/\Q$가 유한 갈루아 확장이라 하자. 다음이 동치임을 증명하라.
\begin{enumerate}[label=(\roman*)]
  \item $E\subseteq\mathcal C_0$.
  \item $\Gal(E/\Q)$가 $2$-군이다.
  \item 모든 중간체를 적절히 정렬하여 $\Q$에서 $E$까지 차수 $2$인 확장들의 탑을 만들 수 있다.
\end{enumerate}
\end{exercise}
